% Round 15. A concrete failure of an open-edge-only extension.
\section{Positive-delta mass at a moving tree vertex}
\label{r15vg:section}

The PWS moving-tree proof cannot be extended to arbitrary actual
inputs by discarding the lower-dimensional vertex strata.  The
following example identifies a precise failed step; it is not a
counterexample to existence of an ordinary projected form.

\begin{proposition}[A vertex contribution missed by open-edge integration]
\label{r15vg:vertex}
Let $K/\mathbf Q_p$ be any finite valued extension.  There are a
global actual form
\[
 \alpha\in B^{2,2}(\mathbf P^1_x\times\mathbf P^1_z)
\]
and a nonzero global PWS form $\omega\in\operatorname{PS}^{1,1}
(\mathbf P^1_x)$ with these properties.
\begin{enumerate}
\item For projection $\pi$ to the first factor, its prescribed
proper current is $\pi_*[\alpha]=[\omega]$, with the original
$K$ normalization.
\item On a compact ordinary base annulus, the relevant marked
fiber tree has marks $0,x,\infty$.  On every strict open edge
of this tree, the restriction of $\alpha$ is zero.  Its nonzero
top measure is carried by the moving branch vertex of this
tree, whose radius is $|x|$.
\item Consequently a procedure that extends the PWS ray formula
by retaining only its open-edge integrals and discarding all
vertex-supported strata returns zero for this input and fails
the prescribed current pairing.
\end{enumerate}
\end{proposition}

\begin{proof}
Let $u=-\log|x|$ on the base torus and choose a nonnegative,
nonzero $F\in C_c^\infty(\mathbf R)$.  The smooth form
\[
                  \omega=F(u)\,d'u\wedge d''u
\tag{VG15.1}\label{r15vg:omega}
\]
extends by zero to the base projective line.  Its support is
contained in a compact annulus away from $0$ and $\infty$,
and its native integral is nonzero.

Let $\theta_0$ be the standard model-metric curvature of
$\mathcal O_{\mathbf P^1_w}(1)$.  Its top measure is the
native degree-one Dirac mass at the Gauss point.  On
$\mathbf G_{m,x}\times\mathbf P^1_z$ the map
\[
          A(x,z)=(x,w=z/x)
\tag{VG15.2}\label{r15vg:automorphism}
\]
is an algebraic automorphism over the base.  Put
\[
       \theta_x=A^*\operatorname{pr}_w^*\theta_0,
       \qquad \alpha=\theta_x\wedge\pi^*\omega
\tag{VG15.3}\label{r15vg:alpha}
\]
there.  Because the base support is contained in the compact
annulus just chosen, this actual form is zero on a whole
ordinary neighborhood of the two omitted base fibers and
extends by zero to the projective product.  This uses a
genuine analytic pullback and a literal ordinary zero
extension, not restriction of a polyhedral current to a
half-open cell.

For every ordinary smooth scalar test $\eta$ on the proper
base, the frozen projective degree formula (R10.1), applied
to the actual top form $\omega\eta$, gives
\[
 \int_{\mathbf P^1\times\mathbf P^1}
                  \alpha\wedge\pi^*\eta
       =\int_{\mathbf P^1}\omega\eta.
\tag{VG15.4}\label{r15vg:pairing}
\]
The same equality follows directly from the genuine product
measure in coordinates $(x,w)$ and the automorphism $A$.
Its degree is one, so no extension-field factor enters.
Taking $\eta=1$ proves that this current is nonzero.

In the $(x,w)$ coordinates, the top measure in
\eqref{r15vg:alpha} is the product of the PWS base measure
and the Gauss Dirac mass.  On the product skeleton its
polyhedral description is therefore a density in the
$u$ direction multiplied by the divisor $v(w)=0$.
Writing $t=v(z)$, this becomes the graph
\[
                              t=u.
\tag{VG15.5}\label{r15vg:moving-vertex}
\]
With the additional genuine unit $z-x$, the fiber tree with
marks $0,x,\infty$ has three strict open edges: the edge
toward zero, the edge toward $x$, and the edge toward
infinity.  Under $w=z/x$ these are the three open edges
of the tree with marks $0,1,\infty$.  The Gauss curvature
vanishes on each of them and is concentrated at their
common vertex.  Thus the same statement holds for the
pulled-back top measure and for its polyhedral actual
restrictions.  Its vertex moves with $u$ as in
\eqref{r15vg:moving-vertex}.

Every open-edge integral is consequently zero, whereas
\eqref{r15vg:pairing} is nonzero for $\eta=1$.
This proves the asserted failure of the specified procedure.
\end{proof}

\paragraph{The local interface isolated by the example.}
For a finite splitting cover and a finite genuine tuple of base
units, one may compute the push of the positive-delta source
current on a refined moving-tree skeleton, retaining vertices
and every higher-codimension stratum.  This supplies a candidate
balanced polyhedral current on the base skeleton.  To obtain an
ordinary actual output one needs a single ambient
$\delta$-form in those base unit coordinates whose genuine
pullbacks give the candidate on every further test-moment
refinement.  The later universal graphic kernel of
Section~\ref{r15ug:section} and the finite ambient atlas of
Section~\ref{r15ad:section} supply precisely this interface.
Theorem~\ref{r15a2:projection} consequently resolves arbitrary
actual $(2,2)$ input over every smooth algebraic base, with
all ordinary smooth compact mixed tests.  The independent
proper-curve argument in Section~\ref{r15tc:section} additionally
gives every actual scalar test in its top-degree scope.

There are two distinct issues.  First, the PWS calculation's
discarded vertex strata now carry nonzero mass, as the
proposition proves.  Second, a balanced current on one tropical
skeleton need not be a $\delta$-form there: Mihatsch,
\emph{$\delta$-Forms on Lubin--Tate spaces}, Example~2.45(3),
explicitly warns that a flat tropical pushforward can be
balanced without admitting the required $\delta$ presentation.
The verified finite trace was not the missing interface: the
output has degree $(1,1)$, already within its actual range.
The new proof supplies ordinary representability before that
descent.  The example continues to rule out discarding vertex
strata as a shortcut to that proof.

\paragraph{The universal-family route and its subsequent completion.}
The split root functions define an algebraic map to a
configuration space of distinct marked points.  Its root
differences give the corresponding valuated tree.  On a chamber
of fixed tree type, the complete balanced tree cycle can be
written as the sum of its translated vertex-star fans minus
one entire affine line for each bounded edge.  Along a bounded
edge, the two rays from its ends overlap exactly in the segment;
subtracting the full line leaves that segment.  Every remaining
unbounded ray occurs once.  This is an identity of weighted
one-dimensional cycles on that chamber.

Its affine translations permit ambient affine $\delta$ calculus
inside the chamber; Section~\ref{r15ft:section} proves this
ordinary actual output for a supplied single ambient input
and a fixed tree type with positive bounded-edge lengths.
Across a face where a bounded edge
collapses, the vertex positions are only piecewise affine in
the parameters.  One cannot pull a positive-delta input through
those maps by the PWS pullback rule.  A universal balanced
family, an ambient representative for its projected currents,
and compatibility with the analytic configuration-map pullback
resolve this specific problem.  They are proved separately
in Sections~\ref{r15ug:section} and~\ref{r15ad:section}, using
the complete-graph matroid family, its ambient cutting kernel,
and genuine ambient cutoff partitions.  This example itself
does not prove those assertions.  Neither a tropical
forgetful-map picture alone nor equality of individual fiber
integrals is used as an actual-output theorem.
